\documentclass{clean_cv}

% Add a BibTeX-style file encoding all of your publications to include here. You can export this from Zotero. Only include
% publications you want to appear here!
\addbibresource{publications.bib}
\usepackage{hyperref}

\author{Tejasvi Kothapalli}


\begin{document}

\maketitle
% In this section, you can use any of the FontAwesome icons. The commands \faCenter and \faCenterStyle have been defined to properly center the icons
% when using the default font settings.
%
% You can use any of the icons listed in the fontawesome5 package documentation (https://ctan.math.utah.edu/ctan/tex-archive/fonts/fontawesome5/doc/fontawesome5.pdf)
% If you need to specify a specific style (as is done here for the address card), you should use the two-argument \faCenterCycle command
\begin{center}
\begin{tabular}{lll}
    \faCenter{envelope} \href{mailto:tejasvi.kothapalli@berkeley.edu}{tejasvi.kothapalli@berkeley.edu}  & \faCenter{phone-alt} 408-802-0896 &
    % \faCenter{globe} \href{https://tejasvikothapalli.github.io/}{tejasvikothapalli.github.io}
\end{tabular}
\end{center}

\vspace{-1.5em}

\section{Education}

% The datetabular environment takes one argument, which is the width of the left date column. As seen here:
%   9em is a good choice for "dual-date" formats (e.g. Sep 2015 - Nov 2019).
%   4em is a good choice for month/year dates (Sep 2014).
%   2em is a good choice for year-only dates (as seen in the publications)
\begin{datetabular}{9em}
% This is just a tabular environment, for the most part. The dateentry command has been defined for
% convienence. It takes two arguments, the first is the date and the second is whatever you wish placed to the right.

\dateentry{Aug 2023 -- Present}{
\textbf{University of California, Berkeley -- Vision Science Ph.D.}

\textit{\textbf{Area:}} Computational Neuroscience \& NeuroAI

\textit{\textbf{Advisor:}} Jacob Yates,
\textit{\textbf{Thesis Committee:}} Bruno Olshausen, Jack Gallant, Alex Huk, Jacob Yates

\textit{\textbf{Labs:}} The Active Vision and Neural Computation Lab (Led by Jacob Yates), The Redwood Center for Theoretical Neuroscience

\textit{\textbf{Cumulative GPA:}} 3.979

\textit{\textbf{Course Work:}}  \textbf{VS 260A}: Optical and Neural Limits to Vision (A),
\textbf{VS 260C}: Visual Neuroscience (A), \textbf{VS 260D}: Seeing in Time, Space and Color (A+),
\textbf{VS 265}: Neural Computation (A)
}


\dateentry{Aug 2018 -- May 2022}{
\textbf{University of California, Berkeley -- Electrical Engineering \& Computer Science B.S.}

\textit{\textbf{Cumulative GPA:}} 3.746

\textit{\textbf{Upper Division Technical GPA:}} 3.909
% 3.916

\textit{\textbf{Upper Division Coursework:}}
 \textbf{Math 110}: Linear Algebra (A-),
 \textbf{EE 120}: Signals and Systems (P), \textbf{EE 126}: Probability and Random Processes (P), \textbf{EE 127}: Optimization Models in Engineering (A-), \textbf{CS 100}: Principles \& Techniques of Data Science (A), \textbf{CS 161}: Computer Security (A), \textbf{CS 170}: Efficient Algorithms and Intractable Problems (A), \textbf{CS 182}: Designing, Visualizing and Understanding Deep Neural Networks (A-), \textbf{CS 188}: Introduction to Artificial Intelligence (A-), \textbf{CS 189}: Introduction to Machine Learning (A),  \textbf{CS 194-26}: Intro to Computer Vision and Computational Photography (A), \textbf{CS 194-80}: Full Stack Deep Learning (A), \textbf{CS H196A}: Senior Honors Thesis Research (A+), \textbf{CS 280}: Computer Vision (A), \textbf{CS 285}: Deep Reinforcement Learning (A), \textbf{CS 288}: Natural Language Processing (A) 

\textit{\textbf{Research:}} Worked with Professor \textbf{Stella Yu}, Professor \textbf{Meng Lin}, Professor Yubei Chen, Professor Peter Wang. Senior Honors Thesis: \textbf{Studying Dry Eye Syndrome with Machine Learning}.
}

% \dateentry{Fall 2022 - Spring 2023}{
% \textbf{University of California, Berkeley -- Electrical Engineering \& Computer Science 5th Year Masters}

% Pursuing research in Computer Vision under Stella Yu. Will apply to PhD programs after completing masters.

% \eatvspace}
\end{datetabular}

\section{Experience}

\begin{datetabular}{9em}

\dateentry{Aug 2023 - Present}{
\textbf{Graduate Student Researcher in \textit{The Active Vision and Neural Computation Lab}}

Studying neural recordings from free-viewing primates. Building a foundation model from primate neurophysiology recordings. Also reconstructing images from neural activity.
% Researching Natural Vision in Primates and its applications to Computer Vision.

\eatvspace}

\dateentry{Jun 2020 - May 2023}{
\textbf{Researcher at ICSI (International Computer Science Institute)}

\textit{\textbf{The Effects of Whitening in Neural Networks:}} We modified Batch Normalization to decorrelate the feature map. We found whitening to improve ResNet training on CIFAR-100 by over 2 percent. Collaborated with Professor Yu, Yubei Chen, and Peter Wang.

\textit{\textbf{Meibography Artifical Intelligence:}} Used computer vision and classical machine learning techniques to predict eye diseases. Collaborated with Professor Yu, Professor Lin, and Peter Wang. 

\eatvspace}
\dateentry{Aug 2021 - May 2023}{
\textbf{Student Research Assistant at CRC (Clinical Research Center)}

\textit{\textbf{EasyTear Lipid Layer Analysis:}} Using computer vision techniques on videos of eye to determine lipid layer motion and thickness. Collaborated with Professor Lin, Professor Yu, Peter Wang.

% \textit{\textbf{Tear Aqueous Production Rate:}} Built clinical tool to help compute tear aqueous production rate. Collaborated with Professor Lin.

\eatvspace}

\dateentry{Aug 2022 - Dec 2022}{
\textbf{Machine Learning Engineer at Aizip}

Startup in the tinyML space where ML models are deployed to IoT devices. Worked on people detection and fall detection.

\eatvspace}



% \dateentry{Jan 2019 - Aug 2019}{
% \textbf{Weight Lifting Posture Checker}

% Created an application to check the form of weightlifting movements and provide feedback to users. Working with fitness industry partners like Starting Strength and Stronglifts 5x5 to deploy to consumers. Worked with numerous computer vision techniques: Data Annotation and Augmentation, Object Tracking, Training Object Detection Models with Turicreate, and Human Pose Detection Models with Tensorflow.

% \eatvspace}

\dateentry{May 2017 - Aug 2017}{
\textbf{NASA Ames Research Center Intern}

Worked in the Tensegrity Robot Division. Contributed to an open source web based tensegrity robot simulator.
Used machine learning evolutionary algorithm to locomote twelve-rod tensegrity structures in simulation.




\eatvspace}
% There is something seriously funky happening between the tabular commands and the end of the itemize blocks.
% The command \eatvspace should be used if extra space appears at the end of a datetabular environment.
\end{datetabular}

\section{Publications}
\begin{datetabular}{9em}

\dateentry{1}{
\textbf{Reconstructing Natural Images from Brain Activity during Natural Vision} \newline
\textbf{Tejasvi Kothapalli}*, Ryan Ressmeyer*, Declan Rowley*, Gregory D. Horwitz, Jude Mitchell, Alexander C. Huk, Bruno Olshausen, Jacob Yates. \textit{arXiv, 2026}

\eatvspace}

\dateentry{2}{
\textbf{Active sensing explains cortical variability and improves visual information} \newline
Ryan Ressmeyer*, Declan Rowley*, \textbf{Tejasvi Kothapalli}*, Dekel Galor, Ruei-Jr Wu, Michele Rucci, Jude Mitchell, Daniel Butts, Alexander C. Huk, Jacob Yates. \textit{bioRxiv, 2026}

\eatvspace}

\dateentry{3}{
\textbf{Artificial Intelligence Models Utilize Lifestyle Factors to Predict Dry Eye-Related Outcomes} \newline
Andrew Graham, Jiayun Wang, \textbf{Tejasvi Kothapalli}, Jennifer Ding, Helen Tasho, Alisa Molina, Vivien Tse, Sarah M. Chang, Stella X. Yu, Meng C. Lin. \textit{Nature Scientific Reports, 2025}

\eatvspace}

\dateentry{4}{
\textbf{A Machine Learning Approach to Predicting Dry Eye-Related Signs, Symptoms and Diagnoses from Meibography Images} \newline
Andrew Graham, \textbf{Tejasvi Kothapalli}, Jiayun Wang, Jennifer Ding, Vivien Tse, Penny A. Asbell, Stella X. Yu, Meng Lin.
\textit{Heliyon, 2024}

\eatvspace}

\dateentry{5}{
\textbf{Tracking the Dynamics of the Tear Film Lipid Layer} \newline
\textbf{Tejasvi Kothapalli}, Charlie Shou, Peter Wang, Tatyana Svitova, Andrew Graham, Meng Lin, Stella Yu. \textit{Workshop at Neural Information Processing Systems (Neurips): Medical Imaging}, 2022

\eatvspace}






\dateentry{6}{
\textbf{Saving Energy in Homes Using Wi-Fi Device Usage Patterns} \newline
\textbf{Tejasvi Kothapalli}.
\textit{International Journal of Energy Optimization and Engineering (IJEOE)}, 2018

\eatvspace}

\dateentry{7}{
\textbf{Controlling Tensegrity Robots through Evolution using Friction based Actuation} \newline
\textbf{Tejasvi Kothapalli}, Adrian Agogino.
\textit{NASA Technical Reports}, 2017

\eatvspace}

% There is something seriously funky happening between the tabular commands and the end of the itemize blocks.
% The command \eatvspace should be used if extra space appears at the end of a datetabular environment.
\end{datetabular}


\section{Abstracts}
\begin{datetabular}{9em}

\dateentry{1}{
\textbf{Active sensing explains cortical variability and improves visual information} \newline
Declan Rowley*, Ryan Ressmeyer*, \textbf{Tejasvi Kothapalli}*, Alexander C. Huk, Jude Mitchell, Daniel Butts, Jacob Yates. \textit{Society for Neuroscience, 2026}

\eatvspace}

\dateentry{2}{
\textbf{Standard models fail to capture positional invariances in marmoset foveal V1} \newline
 \textbf{Tejasvi Kothapalli*}, Ryan Ressmeyer*, Declan Rowley, Jude Mitchell, Jacob Yates. \textit{Society for Neuroscience, 2025}

\eatvspace}

\dateentry{3}{
\textbf{Lifestyle and behaviors: predicting clinical signs and symptoms with machine learning} \newline
Meng C. Lin, Andrew D. Graham, \textbf{Tejasvi Kothapalli}, Peter Wang, Jennifer Ding, Vivien Tse, Stella X. Yu.
\textit{Association for Research in Vision and Ophthalmology (ARVO), 2023}

\eatvspace}

\dateentry{4}{
\textbf{AI provides deeper understanding of Meibomian gland morphology and function} \newline
Meng C. Lin, Peter Wang, \textbf{Tejasvi Kothapalli}, Andrew D. Graham, Stella X. Yu. \textit{American Academy of Optometry (AAOPT), 2022}

\eatvspace}


% There is something seriously funky happening between the tabular commands and the end of the itemize blocks.
% The command \eatvspace should be used if extra space appears at the end of a datetabular environment.
\end{datetabular}

\section{Honors, Fellowships \& Selected Programs}
\begin{datetabular}{9em}




\dateentry{Aug 2026 -- Present}{
\textbf{NIH Training Grant in Vision Science (NEI T32-EY007043)}

Funded for graduate research through the National Eye Institute training grant at UC Berkeley.

\eatvspace}

\dateentry{Aug 2026}{
\textbf{MIT Brains, Minds and Machines Summer Course}

Selected to attend a fully funded intensive course at MBL Woods Hole focused on neuroscience, cognitive science, computer science, and biologically inspired AI.

\eatvspace}

\dateentry{July 2026}{
\textbf{CSHL Computational Neuroscience: Vision Course}

Selected to attend a long-running Cold Spring Harbor Laboratory course in computational visual neuroscience, focused on neural coding, decoding, attention, decision-making, and perception.

\eatvspace}

\dateentry{Aug 2023 - May 2025}{
\textbf{CIVO Fellowship}

Generously funded for graduate studies to conduct research which promotes innovative display, graphics, and optical technology for the healthy and diseased eye.

\eatvspace}

\dateentry{Aug 2021 - May 2022}{
\textbf{EECS Honors Program}

Recognizes EECS students who commit to research, strong academics, and writing a senior thesis.

\eatvspace}



% There is something seriously funky happening between the tabular commands and the end of the itemize blocks.
% The command \eatvspace should be used if extra space appears at the end of a datetabular environment.
\end{datetabular}

\section{Teaching \& Services}
\begin{datetabular}{9em}

\dateentry{Aug 2025 - Dec 2025}{
\textbf{GSI for Neuroscience 172L: Cognitive and Computational Lab}

Course teaches experimental and analytical techniques used by cognitive and computational neuroscientists. Led lab sections, content creation, grading, and office hours.

\eatvspace}

\dateentry{Aug 2023 - September 2024}{
\textbf{Program Committee Chair, Bay Area Vision Research Day (BAVRD)}

Help plan, fundraise, and invite speakers and abstract presenters.

\eatvspace}

\dateentry{Aug 2023 - Present}{
\textbf{CIVO BASIS Coordinator}

I help organize and teach about general Vision Science to 4th graders around the East Bay. We work with the organization Bay Area Scientists Inspiring Students (BASIS).

\eatvspace}

\dateentry{May 2022 - Present}{
\textbf{Teacher at Inspirit AI}

Teaching machine learning to high students at Khan Lab School, Bellarmine College Preparatory, Evergreen Valley High School, and Bentley School.  

\eatvspace}



% There is something seriously funky happening between the tabular commands and the end of the itemize blocks.
% The command \eatvspace should be used if extra space appears at the end of a datetabular environment.
\end{datetabular}

\section{Talks}
\begin{datetabular}{9em}
\dateentry{November 2026}{
\textbf{Redwood Meeting:} From $1/f$ to Diffusion: State-of-the-Art Brain Decoding during Natural Vision

\eatvspace}

\dateentry{December 2025}{
\textbf{CIVO Day:} Measuring Sensitivity of V1 activity to eye movements

\eatvspace}

\dateentry{October 2025}{
\textbf{Vision Science Retreat:} Standard models fail to capture positional invariances in marmoset foveal V1

\eatvspace}

\dateentry{September 2025}{
\textbf{Bay Area Vision Research Day (BAVRD):} Towards a Foundation Model for Primate Visual Cortex

\eatvspace}

\dateentry{May 2025}{
\textbf{Redwood Meeting:} V1 Sensitivity to Eye Movements \textbf{and} Predicting V1 with Sparse Coding

\eatvspace}

\dateentry{December 2024}{
\textbf{CIVO Day:} High-accuracy retinal input tracking in free-viewing primates

\eatvspace}


\end{datetabular}

\section{References}
\begin{datetabular}{9em}
\dateentry{1}{
\textbf{Jacob Yates}, yates@berkeley.edu \newline
Professor, Herbert Wertheim School of  Vision Science and Optometry, UC Berkeley \newline
Director, Redwood Center for Theoretical Neuroscience

\eatvspace}


\dateentry{2}{
\textbf{Bruno Olshausen}, baolshausen@berkeley.edu \newline
Professor, Neuroscience and Herbert Wertheim School of Optometry \& Vision Science, UC Berkeley

\eatvspace}
% \dateentry{2}{
% \textbf{Stella Yu}, stellayu@berkeley.edu \newline
% Professor, Electrical Engineering and Computer Sciences, University of Michigan, Ann Arbor \newline 
% Adjunct Professor, Electrical Engineering and Computer Sciences, UC Berkeley \newline
% Director, ICSI Vision Group

% \eatvspace}

% \dateentry{3}{
% \textbf{Meng Lin}, mlin@berkeley.edu \newline
% Professor, Herbert Wertheim School of  Vision Science and Optometry, UC Berkeley \newline
% Director, Clinical Research Center

% \eatvspace}

% \dateentry{4}{
% \textbf{Yubei Chen}, yubeic@fb.com \newline
% Professor, Electrical and Computer Engineering, UC Davis

% \eatvspace}
\dateentry{3}{
\textbf{Alex Huk}, alexhuk@g.ucla.edu \newline
Professor, Departments of Psychiatry and Ophthalmology, David Geffen
School of Medicine, UCLA
Director, Fuster Laboratory \newline
\eatvspace}

\dateentry{4}{
\textbf{Peter Wang}, peterwg@berkeley.edu \newline
Assistant Professor, Computational Science and Engineering, Georgia Institute of Technology

\eatvspace}

\dateentry{5}{
\textbf{Michael Silver}, masilver@berkeley.edu \newline
Professor, Neuroscience and Herbert Wertheim School of Optometry \& Vision Science, UC Berkeley \newline 
Director, UC Berkeley Center for the Science of Psychedelics

\eatvspace}

\dateentry{6}{
\textbf{Frederic Theunissen}, theunissen@berkeley.edu \newline
Professor, Neuroscience and Integrative Biology, UC Berkeley \newline 
\eatvspace}


% There is something seriously funky happening between the tabular commands and the end of the itemize blocks.
% The command \eatvspace should be used if extra space appears at the end of a datetabular environment.
\end{datetabular}



\end{document}
